%  LaTeX support: latex@mdpi.com 
%  For support, please attach all files needed for compiling as well as the log file, and specify your operating system, LaTeX version, and LaTeX editor.

%=================================================================
\documentclass[algorithms,article,submit,pdftex,moreauthors]{Definitions/mdpi} 

\graphicspath{{figures/}} % path to folder with figures
\usepackage{caption}
\usepackage{subcaption}
%\usepackage[export]{adjustbox}
\usepackage{listings}
\usepackage{xcolor}
\usepackage{graphicx}
\usepackage{gensymb} % for degree symbol
\usepackage{lscape}
\usepackage{pxfonts}

\definecolor{codegreen}{rgb}{0,0.6,0}
\definecolor{codegray}{rgb}{0.5,0.5,0.5}
\definecolor{codepurple}{rgb}{0.58,0,0.82}
\definecolor{backcolour}{RGB}{255,255,255}
\definecolor{SlateGray}{HTML}{708090}
%\usepackage{algorithm}
%\usepackage{algpseudocode}

\lstdefinestyle{mystyle}{
    backgroundcolor=\color{backcolour},   
    keywordstyle=\bfseries,
    numberstyle=\tiny\color{codegray},
    stringstyle=\color{SlateGray},
    basicstyle=\ttfamily\footnotesize,
    breakatwhitespace=false,         
    breaklines=true,                 
    captionpos=t, 
    keepspaces=true,                 
    numbers=left,                    
    numbersep=5pt,                  
    showspaces=false,                
    showstringspaces=false,
    showtabs=false,                  
    tabsize=2,
    frame=topline|bottomline,
    }

\lstset{style=mystyle}
\captionsetup[lstlisting]{position=top, margin=0cm,labelfont={bf, small, stretch=1.17}, labelsep=colon, textfont={small, stretch=1.17}, aboveskip=6pt, singlelinecheck=off, justification=justified}
	
\usepackage{soul} % highlighting text

%--------------------
% Class Options:
%--------------------
%----------
% journal
%----------

%---------
% article
%---------

%=================================================================
% MDPI internal commands
\firstpage{1} 
\makeatletter 
\setcounter{page}{\@firstpage} 
\makeatother
\pubvolume{1}
\issuenum{1}
\articlenumber{0}
\pubyear{2022}
\copyrightyear{2022}
%\externaleditor{Academic Editor: Firstname Lastname}
\datereceived{} 
%\daterevised{} % Only for the journal Acoustics
\dateaccepted{} 
\datepublished{} 
%\datecorrected{} % Corrected papers include a "Corrected: XXX" date in the original paper.
%\dateretracted{} % Corrected papers include a "Retracted: XXX" date in the original paper.
\hreflink{https://doi.org/} % If needed use \linebreak
%\doinum{}
%------------------------------------------------------------------
% The following line should be uncommented if the LaTeX file is uploaded to arXiv.org
%\pdfoutput=1

%=================================================================
% Add packages and commands here. The following packages are loaded in our class file: fontenc, inputenc, calc, indentfirst, fancyhdr, graphicx, epstopdf, lastpage, ifthen, lineno, float, amsmath, setspace, enumitem, mathpazo, booktabs, titlesec, etoolbox, tabto, xcolor, soul, multirow, microtype, tikz, totcount, changepage, attrib, upgreek, cleveref, amsthm, hyphenat, natbib, hyperref, footmisc, url, geometry, newfloat, caption

%=================================================================
% Full title of the paper (Capitalized)
\Title{Satellite Image Processing by Python and R Using Landsat 9 OLI/TIRS and SRTM DEM Data on Côte d'Ivoire, West Africa}

% MDPI internal command: Title for citation in the left column
\TitleCitation{Satellite Image Processing by Python and R Using Landsat 9 OLI/TIRS and SRTM DEM Data on Côte d'Ivoire, West Africa}

% Author Orchid ID: enter ID or remove command
\newcommand{\orcidauthorA}{0000-0002-5759-1089} % Add \orcidA{} behind the author's name
\newcommand{\orcidauthorB}{0000-0002-6461-1551} % Add \orcidB{} behind the author's name

% Authors, for the paper (add full first names)
\Author{Polina Lemenkova $^{1}$\orcidA{} and Olivier Debeir $^{1}\orcidB{}$}
% \dagger,\ddagger
%\longauthorlist{yes}

% MDPI internal command: Authors, for metadata in PDF
\AuthorNames{Polina Lemenkova $^{1}$ and Olivier Debeir $^{1}$}

% MDPI internal command: Authors, for citation in the left column
\AuthorCitation{Lemenkova, P.; Debeir, O.}
% If this is a Chicago style journal: Lastname, Firstname, Firstname Lastname, and Firstname Lastname.

% Affiliations / Addresses (Add [1] after \address if there is only one affiliation.)
\address{%
$^{1}$ \quad Université Libre de Bruxelles, École polytechnique de Bruxelles (EPB, Brussels Faculty of Engineering), Laboratory of Image Synthesis and Analysis, Building L, Campus de Solbosch, ULB -- LISA CP165/57, Avenue Franklin D. Roosevelt 50, B-1050, Brussels, Belgium}

% Contact information of the corresponding author
\corres{Correspondence: polina.lemenkova@ulb.be; Tel.: +32471860459 (P.L.)}

% Abstract
\abstract{Programming algorithms for the terrain analysis and environmental monitoring from satellite images have been successful in dealing with challenging multispectral remote sensing data. In this paper, the problem of deriving information on vegetation coverage and relief parameters is addressed, dealing with the case of Landsat 9 OLI/TIRS and SRTM DEM datasets covering the territory of Côte d'Ivoire, West Africa. The vegetation structure in this region is heterogeneous which reflects the complexity of the terrain structure. Due to the high level of spatial detail, high-resolution satellite images require efficient algorithms of data handling, including image resampling, band composition, statistical processing and map algebra for rapid calculation of raster cells. This requires the use of advanced programming methods targeted at image processing. To this end, we used the EarthPy library of Python and 'raster' and 'terra' packages of R for image processing. The methodology included creating colour composites, computing vegetation indices and geomorphic terrain modelling. Compared to the state-of-the-art GIS solutions, the functionality of scripts proved effective handling of spatial data. The scripting algorithms of Python and R demonstrated high quality, speed and effectiveness of code, compared to the existing GIS approaches for remote sensing data processing. The upscaled mapping was shown at the two hierarchical levels: regional level of Yamoussoukro surroundings for environmental analysis, and local level of Kossou Lake for topographic modelling.}

% Keywords
\keyword{Image processing, remote sensing; programming; Python; R; satellite image; Digital Terrain Model; cartography; mapping; spatial analysis} 

% The fields PACS, MSC, and JEL may be left empty or commented out if not applicable
%\PACS{J0101}
%\MSC{}
%\JEL{}

%%%%%%%%%%%%%%%%%%%%%%%%%%%%%%%%%%%%%%%%%%
\begin{document}

\section{Introduction}
Repeated structures representing similar objects on the Earth's surface are commonly depicted on the satellite images, represented by pixels comprising the matrix of the image scene with comparable range of values. Similar values of raster cells are prevailing in the regions with identical land cover types which comprise the mosaic of the diverse land structure on the Earth. Identifying and recognising these feature objects on the surface is possible using the parameters of pixels in the satellite images: brightness, resolution \cite{4426419,322052,5446016}, intensity \cite{4761020,8519176} and texture \cite{995559,899813,8868688}. Algorithms of image classification and processing enable to identify pixel values using binary data derived from raster cells. Such information derived from the interpretation of the satellite images have been traditionally used for the thematic mapping: land cover fractions \cite{COULIBALY2021108092}, vegetation types \cite{7729337,Lemenkova2020d,6946534,Lemenkova202019,6352694}, hydrologic and geomorphic modelling \cite{2015hydrologie}, urban sprawl, geological exploration and mining \cite{1457509,NGOM2022102873}. Besides the satellite images, another source of information for mapping is provided by airborne aeromagnetic surveys \cite{JESSELL2015440,SIAGNE2022104680} and LiDAR data for 3D terrain analysis \cite{SHARMA2021317,8978301,9213897,YEH2017236,9553919,LI201955,GUPTA2022100817}. 

Since remote sensing data present the important source of spatial information, accurate processing of satellite images is essential for the environmental modelling and mapping. The algorithms of image classification and processing are diverse with general idea based on reading the information from the pixels comprising the image. The analysis of this information enables mapping the diverse phenomena, processes and objects based on image classification and interpretation of information derived from Earth observation data \cite{PERIN2021106694,PAHLEVAN2022112860,Lemenkova202141,ABU2021107863,ATTOUMANE2022103285}. The main objective of image processing is deriving this information from the raw data in raster format. In this context, the search for effective methods of image processing resulted in the variety of existing GIS software with specially designed embedded functions for raster data processing. However, in many practical situations, GIS present limitations either due to the commercial licence of software (e.g., ArcGIS, Erdas Imagine, eCognition) or functional restrictions (e.g., SAGA GIS, QGIS). Accordingly, open source programming tool present effective solutions that are able to process satellite imagery using advanced methods of machine learning \cite{MASOLELE2021112600,9066075,MARTINS2022113203,rs71114876}. Recent studies revealed that scripting techniques enable to handle Earth observation data effectively and accurately \cite{6344536,ELLEFSEN20191171,Lemenkova202212,6910592,Lemenkova202215,4160259}. A variety of existing programming languages raises the questions of optimal and effective tools suitable specifically for image processing. 

Python is undoubtedly one of the most successful high-level programming languages \cite{van1995python}. It become one of the most powerful tools that can be used for data science. Besides general applications in data science and machine learning \cite{8757088,DebeirDecaestecker,Debeir2008,9719489}, Python can be applied for modelling Earth phenomena using libraries specially designed for spatial analysis \cite{Ahmed}. Image and signal processing by object-oriented algorithms of Python is pervasive in Earth sciences and provides solutions to a wide class of problems in a variety of domains. Particular examples of using Python have been reported in geodetic studies \cite{Zhou2021GPS2spaceAO}, topographic analysis \cite{9254627}, hydrogeological modelling \cite{Mogaji}, modelling air pollution and land use \cite{Maetal2020}, batch spatial data processing \cite{6057428}, creating panorama images \cite{JASKOLKA2019e02560}, scanning tunneling microscopy \cite{BASTIDAS2021150821}, to mention a few. The general idea behind the image processing is that vegetation properties, topographic relief and land cover types can be depicted using data derived from the cells of the raster images. These data can be interpreted and visualized using the properties of DNs of the pixels which depend on the texture, composition and fabric of the objects shaping the face of the Earth \cite{YUAN2022106966}. 

R language is another fundamental tool for spatial data processing \cite{RCoreTeam}, which also accounts for different statistical components available, a variety of its packages and rich graphical functionality \cite{Murrell}. R and its packages has long been used for data retrieval, analysis, modelling and visualization in various industries including robotics, engineering and GIS \cite{7847824,Lemenkova2022a,9378399,Lemenkova202216,9390011}. R demonstrates superiority in statistical computing and graphics over common GUI-based software for data analysis. Specifically, it excels in big data analysis \cite{7570960,9752002,6597171}, data mining \cite{7978404} and visualization functionality \cite{Lemenkova201991} which includes modelling and plotting supported through the variety of built-in packages. Owing to its long history of stable development, R has a more mature support of rich collection of libraries compared to Python, which contains a wide variety of the mathematical statistical tools, starting from the very basic functions to highly sophisticated algorithms \cite{9675988}. 

Recent advances in geospatial data processing \cite{8226275} and image classification \cite{9844952,8117873} have made R successful for remote sensing data processing including classification and segmentation of complex scenes. This poses us a new way to process and model satellite images and derive spatially structured information from the scenes. By modelling the spatial variations in the landscape structure, such as vegetation and topography, methods of geospatial R libraries 'terra' and 'raster' consider spatial dimensions of the image when manipulating the remote sensing data. Thus, using the similarity between the values of pixels, the feature extraction is defined for the classification of variables, e.g., vegetation types on the scene. In this regard, optimal selection of spectral bands for classifications of the image is indispensable for scene analysis, which is possible in 'terra' package of R. 

An essential advantage of the programming languages for image processing consists in their compatibility. Thus, in contrast to the stand along state-of-the-art GIS software, scripting libraries can be integrated with other GIS via modules, plugins and toolkits \cite{6407912,7378577,VOS2019104528}. For instance, novel libraries of Python enable to perform various tasks of image processing and spatial analysis through integration \cite{1247428,Rey2010PySALAP,Rey2010}. Selected recent examples include, for instance, image interpolation, calibration and correlation \cite{DEDEUSFILHO2022101131}, compressed sensing and image reconstruction \cite{FARRENS2020100402}, identifying and extracting the vertical features \cite{GAO2022105503}, meteorological modelling \cite{6528254} with resulted images in compatible formats. A collection of Python libraries \cite{ROBERTS2022105192,CHEN2022105362} has recently built up with EarthPy, tailored for spatial tasks of geo-information data processing \cite{wasser2019EarthPyAP}. 

Following the idea of using the advanced programming methods used to visualise a representation of object features and classes on the Earth's surface, the goal of this paper is to present the application of Python and R algorithms for image processing. Specific objectives include the following tasks: \emph{i)} plotting the colour composite of the Landsat OLI/TIRS  image by the combinations of three channel sensor bands to display the image; \emph{ii)}; computing the selected vegetation indices using the principles of map algebra based on DNs of reflectance value of pixels in selected spectral bands -- NIR and red \emph{iii)} 
terrain modelling with geomorphic visualization of slope, aspect, hillshade and elevation for the relief analysis using the SRTM DEM data. 

%\newpage
\section{Study Area}

The study ares is focused on Côte d'Ivoire, West Africa, Figure \ref{fig01}. 

\begin{figure}[H]
\centering
	\includegraphics[width=12.0 cm]{Fig_01}
	\caption{Topographic map of Côte d'Ivoire. Mapping software: Generic Mapping Tools (GMT) scripting toolset. Data source: GEBCO/SRTM. Cartography: authors.
	\label{fig01}}
\end{figure}  

The terrain structure here forms the two parts: a coastal part formed by the alluvium and marine sands and a continental part with hilly relief \cite{Thomasset}. The land cover types are associated with the two major ecosystem types -- forest and savanna. The forests are distributed in the southern part of the country, while savanna -- in the north off Bouak\'{e} \cite{Rougerie}. The vegetation types are largely controlled by the equatorial or subequatorial rainfall regime: long dry season is typical for the northern half of Côte d'Ivoire with distributed savannah, while tropical regime characterises southern part with dominating forests \cite{Tricart}. Favourable climatic and environmental setting maintain farming and agricultural activities which make Côte d'Ivoire the world's largest cocoa producer \cite{Losch,Planhol,Smithetal2014,Laderach} and active exporter of food-producing crops \cite{Sawadogo}, such as coffee, bananas, pineapples and palm oil \cite{Pelissier}. 

Besides agriculture, the activities of population include exploration and mining \cite{SAKO2018297,SAUERWEIN2020104323,MURRAY2019351} and limited fishing along the narrow continental shelf \cite{Sournia}. Recently, Côte d'Ivoire experiences environmental changes common for humid West Africa, which include deforestation and fragmentation of ecosystems \cite{Chatelainetal2010,Chatelain}, fluctuations in species richness in savannah biomes \cite{Kolongo}, and increase in habitat heterogeneity which affects biodiversity and species community structure \cite{Yeo}. Besides, urban sprawl is exceptional for Côte d'Ivoire since 1950s, due to the natural demographic growth and acceleration of population in the industrial agglomerations \cite{Dubresson,Chaleard,Cotten1974,Cotten1973}. Geographically, this results in increased areas of cities and modified transport network \cite{Cotten1985}. The surroundings of Kossou Lake were studied as an enlarged fragment of the terrain. This lake has been created artificially following the construction of the dam across the Bandama River, and now forms one of the important inland water bodies of Côte d'Ivoire. Due to its functional environmental aspects, e.g., habitats for dominant biotas and supporting food webs, it plays an essential ecological role in the environment of sub-Saharan Africa \cite{MSISKA2009487}. 

%%%%%%%%%%%%%%%%%%%%%%%%%%%%%%%%%%%%%%%%%%
%\newpage
\section{Materials and Methods}

\subsection{Data}
The data were collected from the USGS GloVis, a satellite image repository with open access to remotely sensed data, Figure \ref{fig02}. 

\begin{figure}[H]
     \begin{subfigure}[r]{0.60\textwidth}
         \centering
         \includegraphics[width=11cm]{Fig_02a}
         \caption{Collecting Landsat 9 OLI/TIRS C2 L1.}
         \label{fig:00a}
     \end{subfigure} \vfill \vspace{2mm}
     \begin{subfigure}[r]{0.60\textwidth}
         \includegraphics[width=11cm]{Fig_02b}
         \caption{Collecting SRTM Void Filled DEM.}
         \label{fig:00a}
     \end{subfigure}
        \caption{Data capture from the USGS Global Visualization Viewer (GloVis) satellite image repository to accessing remote sensing data: Landsat 9 OLI/TIRS C2 L1 and SRTM DEM for terrain analysis}
        \label{fig02}
\end{figure}

The satellite image used for data processing is Landsat 9 OLI/TIRS C2 L1 scene covering the central part of Côte d'Ivoire, \emph{LC09\_L1TP\_197055\_20220111\_20220122\_02\_T1}; Scene Identifier: LC91970552022011LGN01; Land Cloud Coverage: 1.29; Date acquired: 2022/01/11. The space-based images were obtained from the Landsat 9 Earth observation satellite which has two remote sensing instruments onboard -- the Operational Land Imager (OLI-2) and Thermal Infrared Sensor (TIRS-2) for optical and thermals sensors, and nine spectral bands \cite{USGS,Landsat9}. More technical information is available online: \href{https://www.usgs.gov/landsat-missions/landsat-9}{Landsat 9}. The data used for terrain analysis is SRTM Void Filled DEM, from the \href{https://lpdaac.usgs.gov/products/srtmgl3v003/}{SRTM3 NASA/NGA} with 3 arc-second resolution, 90 m. Entity ID: SRTM3N07W006V2; tile coordinates: 5°-6°W, 7°-8°N. Acquisition: 2000-02-11. The general map showing study area has been plotted based on the GEBCO \cite{GEBCO} using methods of GMT scripting \cite{Wessel} described earlier \cite{Lemenkova2022b,Lemenkova202217}. 

%\newpage
\subsection{Image processing in R}
The R packages 'terra' \cite{RTerra} and 'raster' \cite{Hijmans} were used for computing vegetation indices and image processing. The colour scheme resource of RColorBrewer package were used for visualization and plotting \cite{Neuwirth}.

\subsubsection{Deriving information on satellite image properties and spatial statistics}
To cope with a series of Landsat bands and for searching metadata, we examined the properties of the digital images using the regular expression by the command-line utilities of Unix as a combination of grep, gdal and echo, Listing \ref{lst:01}. 

\begin{lstlisting}[language=bash, caption=Unix shell script by GDAL and echo-grep utilities to inspect the geometry of Landsat TIFs, label={lst:01}]
for f in *tif ; do  echo $f; gdalinfo $f | grep Size; done

\end{lstlisting}The spatial resolution described by the pixel size is of 30 m for Bands 1 to 7, 9 and 10, while 15 m for Band 8 (Panchromatic). Figure \ref{fig:03a}. 

\begin{figure}[H]
     \centering
     \begin{subfigure}[b]{0.45\textwidth}
         \centering
         \includegraphics[width=7cm]{Fig_03a}
         \caption{Using grep and echo Unix utilities with 'gdalinfo' module of GDAL to scan metadata}
         \label{fig:03a}
     \end{subfigure} \hfill \hspace{2mm}
     \begin{subfigure}[b]{0.45\textwidth}
         \centering
         \includegraphics[width=7cm]{Fig_03b}
         \caption{Creating SpatRaster objects in R and check spatial data: projection, datum, ellipsoid, etc.}
         \label{fig:03b}
     \end{subfigure}
        \caption{Inspecting the geometry of Landsat bands (a) which shows 30 pixels resolution for all the bands except for Band 8 (panchromatic) with 15-pixels resolution. Checking image properties in R package 'terra' (b) for single Landsat-8 OLI/TIRS C1 bands (1:11)}
        \label{fig03}
\end{figure}

The spatial information of the Landsat channels (bands) is presented in metadata shown in Figure \ref{fig:03b}, inspected by R. It shows the number and location of the spectral measurements, spatial extent, geodetic information (coordinate reference system, ellipsoid and datum) obtained by the GPS of the spacecraft, spatial and the radiometric resolution of the image. Each image (the band of the Landsat scene) is read-in to R SpatRaster object class, which represents the equalised area of rectangular cells as regards the units of the coordinate reference system, Figure \ref{fig:03b}. 

\subsubsection{Visualising single bands and generating composite maps}

The properties of the image were accessed from a SpatRaster object using R commands, Listing \ref{lst:02}, where we inspected the coordinate reference system (CRS), the geometry of the image (Number of cells, rows and columns), bands extent, resolution and origin. The panchromatic channel (Band 8) has a higher resolution (15 m) compared to the set of the multispectral bands in visible and thermal bands of spectra (30 m). To ignore this difference in the matrix of the image, the raster in Band 8 (panchromatic) was resampled to the structure of the same number of rows and columns as in other band (here: target Band 1, Ultra Blue) by created a RasterBrick multi-layer object, Listing \ref{lst:02}. 

\begin{lstlisting}[language=r, caption=R code used for deriving satellite mage information and statistics and inspecting the Landsat 8-9 OLI bands, label={lst:02}]
library(terra)
## terra version 1.2.11
# Change working directory
setwd("/Users/polinalemenkova/Documents/R/Image_Processing/LC09_L1TP_197055_20220111_20220122_02_T1")
b2 <- rast('LC09_L1TP_197055_20220111_20220122_02_T1_B2.TIF') # Blue
b3 <- rast('LC09_L1TP_197055_20220111_20220122_02_T1_B3.TIF') # Green
b4 <- rast('LC09_L1TP_197055_20220111_20220122_02_T1_B4.TIF') # Red
b5 <- rast('LC09_L1TP_197055_20220111_20220122_02_T1_B5.TIF') # Near Infrared (NIR)
# Checking metadata by printing the variables
b2
# Checking coordinate reference system (CRS)
crs(b2)
# Number of cells
ncell(b2)
# Number of rows & columns
dim(b2)
# Image information and statistics spatial resolution
res(b2)
# Number of bands
nlyr(b2)
# Comparing bands extent, number of rows and columns, projection, resolution and origin
compareGeom(b5, b4)
# Creating a SpatRaster with multiple layers from the existing SpatRaster (single layer) objects.
s <- c(b5, b4, b3)
# Checking properties of the multi-band image
s
# Creating a list of raster layers to use
filenames <- paste0('LC09_L1TP_197055_20220111_20220122_02_T1_B', 1:11, ".tif")
filenames
# Resampling raster in Band 8 (panchromatic) to the structure of Band 1:
library(raster)
r1 = raster::brick("LC09_L1TP_197055_20220111_20220122_02_T1_B8.tif")
r2 = raster::brick("LC09_L1TP_197055_20220111_20220122_02_T1_B1.tif")
r1 <- resample(r1, r2, method="bilinear")
s <- stack(r1, r2)
# Creating the multi-layer SpatRaster using the filenames
landsat <- rast(filenames)
landsat
\end{lstlisting} 

The layers were then read in the 'landsat' object where the layers are represented by the bands of Landsat-9 OLI with reflection intensity in the following wavelengths: Ultra Blue, a.k.a. coastal aerosol (Band 1), Blue (Band 2), Green (Band 3), Red (Band 4), Near Infrared (NIR) (Band 5), Shortwave Infrared (SWIR) 1 (Band 6), Shortwave Infrared (SWIR) 2 (Band 7), Panchromatic (Band 8), Cirrus (Band 9), Thermal Infrared (TIRS) 1 (Band 10), Thermal Infrared (TIRS) 2 (Band 11). The visuallisation of the single bands has been performed using Listing \ref{lst:03} and presented as a subplot of grayscale scenes in Figure \ref{fig04}. The visual inspection of the satellite image product built from the monochrome image data (Figure \ref{fig04}) is useful to extract a general information on general large-scale features, such as river systems, lakes, mountains, fracture patterns, large geological lineaments, city extent with notable objects (highways). These objects are visible and can be identified by their shape and size, and interpreted using their spatial properties related to the pixel brightness values. 

\begin{lstlisting}[language=r, caption=R code for plotting the 11 original individual layers (raw bands) of the multi-spectral image Landsat-9 as grayscale bitmap images, label={lst:03}]
library(terra)
## terra version 1.2.11
# Change working directory
setwd("/Users/polinalemenkova/Documents/R/Image_Processing/LC09_L1TP_197055_20220111_20220122_02_T1")
b1 <- rast('LC09_L1TP_197055_20220111_20220122_02_T1_B1.TIF') # Ultra Blue
b2 <- rast('LC09_L1TP_197055_20220111_20220122_02_T1_B2.TIF') # Blue
b3 <- rast('LC09_L1TP_197055_20220111_20220122_02_T1_B3.TIF') # Green
b4 <- rast('LC09_L1TP_197055_20220111_20220122_02_T1_B4.TIF') # Red
b5 <- rast('LC09_L1TP_197055_20220111_20220122_02_T1_B5.TIF') # NIR
b6 <- rast('LC09_L1TP_197055_20220111_20220122_02_T1_B6.TIF') # Green
b7 <- rast('LC09_L1TP_197055_20220111_20220122_02_T1_B7.TIF') # Red
b8 <- rast('LC09_L1TP_197055_20220111_20220122_02_T1_B8.TIF') # Panchrom.
b9 <- rast('LC09_L1TP_197055_20220111_20220122_02_T1_B9.TIF') # Cirrus
b10 <- rast('LC09_L1TP_197055_20220111_20220122_02_T1_B10.TIF') # TIRS 1
b11 <- rast('LC09_L1TP_197055_20220111_20220122_02_T1_B11.TIF') # TIRS 2.
# Plotting natural color composite from bands B4, B3 and B2.
landsatRGB <- c(b4, b3, b2)
plotRGB(landsatRGB, stretch = "lin")
# Plotting 11 original individual layers as raw bands of the Landsat-9 OLI.
par(mfrow = c(4, 3))
plot(b1, main = "Ultra Blue (Band 1)", col = gray(0:100 / 100), legend=FALSE)
plot(b2, main = "Blue (Band 2)", col = gray(0:100 / 100), legend=FALSE)
plot(b3, main = "Green (Band 3)", col = gray(0:100 / 100), legend=FALSE)
plot(b4, main = "Red (Band 4)", col = gray(0:100 / 100), legend=FALSE)
plot(b5, main = "NIR (Band 5)", col = gray(0:100 / 100), legend=FALSE)
plot(b6, main = "SWIR (Band 6)", col = gray(0:100 / 100), legend=FALSE)
plot(b7, main = "NIR (Band 7)", col = gray(0:100 / 100), legend=FALSE)
plot(b8, main = "Panchrom. (Band 8)", col = gray(0:100 / 100), legend=FALSE)
plot(b9, main = "Cirrus (Band 9)", col = gray(0:100 / 100), legend=FALSE)
plot(b10, main = "TIRS 1 (Band 10)", col = gray(0:100 / 100), legend=FALSE)
plot(b11, main = "TIRS 2 (Band 11)", col = gray(0:100 / 100), legend=TRUE)
landsatRGB <- c(b4, b3, b2)
plotRGB(landsatRGB, stretch = "lin")
\end{lstlisting} 

The difference in shades of grey correspond to the numerical values and range of pixel values in the different channels of Landsat, which is owing to the spectral refrectale that differs in various bands for the objects of the Earth surface. Due to the individual reflectance of the solar radiation by features, the values differ accordingly. THis fundamental properties fo features enables to use combinations of various bands for a better represetning diverse phenomena and objects on the earth's surface. For example, particular applications of Landsat bands have for geology (SWIR Band 7, NIR Band 5 and Red Band 4) \cite{VOLESKY2003235}, mineralogical identification of soil and rock (SWIR Band 7, Red Band 4 and Green Band 2), identifying difference between soil and vegetation (Green in Band 3) \cite{Richards,CampbellWynne}.

Subsetting the large files, such as multispectral Landsat dataset was done using the SpatRaster object and 'subset' function for processing only the necessary bands. In this, it enables to ignore the Cirrus, TIRS and other not relevant bands and leaving only the visible spectra. In this way, a a subset of target layers was selected from a SpatRaster and the other bands were removed from the dataset, respectively, using Listing \ref{lst:04}.

\begin{figure}[H]
\makebox[\textwidth][r]{\includegraphics[width=1.3\textwidth]{Fig_04.jpg}}%
	\caption{Original individual 11 layers (raw bands) of the multi-spectral image Landsat-9 OLI/TIRS C1 'LC09\_L1TP\_197055\_20220111\_20220122\_02\_T1' in grayscale visualization: Ultra Blue (Band 1), Blue (Band 2), Green (Band 3), Red (Band 4), Near Infrared (NIR) (Band 5), Shortwave Infrared (SWIR) 1 (Band 6), Shortwave Infrared (SWIR) 2 (Band 7), Panchromatic (Band 8), Cirrus (Band 9), Thermal Infrared (TIRS) 1 (Band 10), Thermal Infrared (TIRS) 2 (Band 11). Plotting: R.}
	\label{fig04}
\end{figure} 

\begin{lstlisting}[language=r, caption=R code for Subsetting large files using SpatRaster object, label={lst:04}]
library(terra)
setwd("/Users/polinalemenkova/Documents/R/Image_Processing/LC09_L1TP_197055_20220111_20220122_02_T1")
filenames <- paste0('LC09_L1TP_197055_20220111_20220122_02_T1_B', 1:11, ".tif")
filenames
landsat <- rast(filenames)
landsat
# Removing the unnecessary four bands: Panchromatic, Cirrus, TIRS1, TIRS2.
landsatsub <- subset(landsat, 1:7)
names(landsatsub)
# Checking Nnumber of bands in the original and new data
nlyr(landsatsub)
# Renaming the bands
names(landsatsub) <- c("ultra-blue", "blue", "green", "red", "NIR", "SWIR1", "SWIR2")
names(landsatsub)
\end{lstlisting} 

To compare bands in the pairwise combinations (Bands 1 vs 2 and Band 4 vs 5) in selected parts of the spectrum, the images were plotted pixel-by-pixel against those from corresponding bands, Listing \ref{lst:05}. The resulting plot shows a narrow correlation ellipse a histogram representing a strong positive correlation between these two bands - Green and Blue, and Red to NIR, Figure \ref{fig05}.

\begin{lstlisting}[language=r, caption=R code for inspecting the correlation between selected Landsat bands, label={lst:05}]
# Relation between bands
library(terra)
pairs(landsatsub[[1:2]], hist=TRUE, cor=TRUE, use="pairwise.complete.obs", maxcells=100000, main = "Correlation of Band 1 (Ultra-blue) against Blue (Band 2)")
pairs(landsatsub[[4:5]], hist=TRUE, cor=TRUE, maxcells=100000, main = "Correlation of Band 4 (Red) against NIR (Band 5)")
\end{lstlisting} 

Cells with low spectral reflectance in a target band mostly have low reflectance in all the rest of the bands, which is caused by the albedo of the Earth's surface. However, the correlation between the Bands 1 and 2 shows more similarity between the bands, since all the pixels almost form the same line, while NIR and Red Bands demonstrate more differences in spectral reflectance of selected objects, which results in the scattered cloud of points representing the pixels on the scenes. The brightness of pixels in each layer indicate the amount of  incident solar radiation, which is reflected by the Earth's surface in a given wavelength diapason, which affect the similarity and difference in spectral bands. 

\begin{figure}[H]
     \centering
     \begin{subfigure}[b]{0.45\textwidth}
         \centering
         \includegraphics[width=7cm]{Fig_05a.jpg}
         \caption{Band 1 against Band 2}
         \label{fig:00a}
     \end{subfigure} \hfill \hspace{2mm}
     \begin{subfigure}[b]{0.45\textwidth}
         \centering
         \includegraphics[width=7cm]{Fig_05b.jpg}
         \caption{Band 4 against Band 5.}
         \label{fig:00a}
     \end{subfigure}
        \caption{Correlation between Landsat bands for selected Landsat-8 OLI/TIRS C1 bands in multispectral image \emph{LC09\_L1TP\_197055\_20220111\_20220122\_02\_T1\_B}.}
        \label{fig05}
\end{figure}

Since there is not much information in the monochrome bands, they were combined into the several composites to create more meaningful plots. Plotting the colour composites of the Landsat OLI/TIRS image is based on the combinations of the three channel sensor bands used as elements of matrix scene. The three bands assigned to each of the three colour primaries (RGB) displayed the image, Listing \ref{lst:06}.

\begin{lstlisting}[language=r, caption=R code used for plotting color composites for bands of the satellite image, label={lst:06}]
library(raster)
blue = raster("LC09_L1TP_197055_20220111_20220122_02_T1_B2.TIF")
green = raster("LC09_L1TP_197055_20220111_20220122_02_T1_B3.TIF")
red = raster("LC09_L1TP_197055_20220111_20220122_02_T1_B4.TIF")
rgb = stack(red, green, blue)
plotRGB(rgb, scale=65535)
plotRGB(rgb, stretch="lin")
plotRGB(rgb, stretch="hist")
# False color composites
# Color Infrared: colors near-infrared (band 5) as red, red (band 4) as green, and green (band 3) as blue
green = raster("LC09_L1TP_197055_20220111_20220122_02_T1_B3.TIF")
red = raster("LC09_L1TP_197055_20220111_20220122_02_T1_B4.TIF")
near.infrared = raster("LC09_L1TP_197055_20220111_20220122_02_T1_B5.TIF")
rgb = stack(near.infrared, red, green)
plotRGB(rgb, scale=65535, stretch="lin")
# Landsat 9 OLI/TIRS colour composite image, RGB=753
green = raster("LC09_L1TP_197055_20220111_20220122_02_T1_B7.TIF")
red = raster("LC09_L1TP_197055_20220111_20220122_02_T1_B5.TIF")
near.infrared = raster("LC09_L1TP_197055_20220111_20220122_02_T1_B3.TIF")
rgb = stack(near.infrared, red, green)
plotRGB(rgb, scale=65535, stretch="hist")
# Landsat 9 OLI/TIRS colour composite image, RGB=765.
red = raster("LC09_L1TP_197055_20220111_20220122_02_T1_B7.TIF")
green = raster("LC09_L1TP_197055_20220111_20220122_02_T1_B6.TIF")
blue = raster("LC09_L1TP_197055_20220111_20220122_02_T1_B5.TIF")
rgb = stack(red, green, blue)
plotRGB(rgb, scale=65535, stretch="hist")
# Landsat 9 OLI/TIRS colour composite image, RGB=742.
red = raster("LC09_L1TP_197055_20220111_20220122_02_T1_B7.TIF")
green = raster("LC09_L1TP_197055_20220111_20220122_02_T1_B4.TIF")
blue = raster("LC09_L1TP_197055_20220111_20220122_02_T1_B2.TIF")
rgb = stack(red, green, blue)
plotRGB(rgb, scale=65535, stretch="hist")
# Landsat 9 OLI/TIRS colour composite image, RGB=573.
red = raster("LC09_L1TP_197055_20220111_20220122_02_T1_B5.TIF")
green = raster("LC09_L1TP_197055_20220111_20220122_02_T1_B7.TIF")
blue = raster("LC09_L1TP_197055_20220111_20220122_02_T1_B3.TIF")
rgb = stack(red, green, blue)
plotRGB(rgb, scale=65535, stretch="hist")
# Landsat 9 OLI/TIRS colour composite image, Color Infrared: RGB=543
green = raster("LC09_L1TP_197055_20220111_20220122_02_T1_B5.TIF")
red = raster("LC09_L1TP_197055_20220111_20220122_02_T1_B4.TIF")
near.infrared = raster("LC09_L1TP_197055_20220111_20220122_02_T1_B3.TIF")
rgb = stack(near.infrared, red, green)
plotRGB(rgb, scale=65535, stretch="hist")
# Landsat 9 OLI/TIRS colour composite image, Color Infrared: RGB=654
green = raster("LC09_L1TP_197055_20220111_20220122_02_T1_B6.TIF")
red = raster("LC09_L1TP_197055_20220111_20220122_02_T1_B5.TIF")
near.infrared = raster("LC09_L1TP_197055_20220111_20220122_02_T1_B4.TIF")
rgb = stack(near.infrared, red, green)
plotRGB(rgb, scale=65535, stretch="hist")
\end{lstlisting}

The vegetation indices were computed based on the arithmetically transformed combination of bands. This results from the existing formulae for band computations using mathematical operations among the cell brightness in the scenes using addition, subtraction and division of the brightnesses of the NIR and Red Bands of Landsat image. Differences in bands highlight the regions of enhanced and more healthy vegetation in the target area of image which is useful for environmental monitoring and ecological assessment. The resultant image shows the brightness values of vegetation modified by colour palettes made using arithmetic operations on composite bands. 

The most widely used and well known vegetation index is the Normalised Difference Vegetation Index (NDVI) which applies the ratio of the difference in values between Red and NIR bands and their sum. Similar to the NDVI, Enhanced Vegetation Index (EVI) quantifies vegetation greenness with added corrections for atmospheric conditions and canopy background noise. Besides, EVI is more sensitive in forest areas with more dense canopy and vegetation coverage, which is useful for the forest regions of Côte d'Ivoire. The Soil Adjusted Vegetation Index (SAVI) is built up using NDVI with correction for the effects from soil brightness in the regions where vegetative cover is low, e.g. newly deforested areas, desert and fragmented savannah of Côte d'Ivoire. 

Computing vegetation indices is an important environmental tool based on the information from pixel value, which in turn is based on spectral reflectance of Red and NIR bands measured by the sensor. It is commonly computed with a spatial statistics and algebra of bands, using cell values defined on the scene, Listing \ref{lst:07}. The reason for selecting NIR and Red for computing the vegetation indices is because the vegetation appears brighter in these bands since it reflects more energy in NIR/Red bands compared to other wavelengths. Contrarily, water absorbs most of the energy in the NIR band and is represented by black to very dark brown colours. Through a series of interactions among scene elements, they are allocated to the classes based on the weighted values of their relative proportions.

\begin{lstlisting}[language=r, caption=R code used for computing vegetation indices from bands of the satellite image, label={lst:07}]
# Computing vegetation indices
library(terra)
library(RColorBrewer)
setwd("/Users/polinalemenkova/Documents/R/Image_Processing/LC09_L1TP_197055_20220111_20220122_02_T1")
# 1. Normalized Difference Vegetation Index (NDVI) = (NIR - R) / (NIR + R) NDVI = (Band 5 - Band 4) / (Band 5 + Band 4).
vi <- function(img, k, i) {
  bk <- img[[k]]
  bi <- img[[i]]
  vi <- (bk - bi) / (bk + bi)
  return(vi)
}
# For Landsat NIR = 5, red = 4.
ndvi <- vi(landsat, 5, 4)
plot(ndvi, col=rev(terrain.colors(20)), font.main = 1, main = "NDVI for Landsat-9 OLI/TIRS C1 image LC09_L1TP_197055_20220111_20220122_02_T1: Yamassoukro, central C\^ote d'Ivoire", cex.main=0.9)
# Plotting histogram of the NDVI
hist(ndvi, font.main = 1, main = "NDVI values for Landsat-9 OLI/TIRS C1 image \n LC09_L1TP_197055_20220111_20220122_02_T1, Yamassoukro, central C\^ote d'Ivoire", xlab = "NDVI", ylab= "Frequency",
    col = "darkolivegreen1", xlim = c(-0.5, 1),  breaks = 30, xaxt = "n")
axis(side=1, at = seq(-0.6, 1, 0.1), labels = seq(-0.6, 1, 0.1))
#
# 2. Soil Adjusted Vegetation Index (SAVI) corrects NDVI for the influence of soil brightness in areas where vegetative cover is low
# SAVI = ((Band 5 - Band 4) / (Band 5 + Band 4 + 0.5)) * (1.5).
savi_fun = function(nir, red){
  ((nir - red) / (nir + red + 0.5)) * 1.5
}
savi = lapp(landsat[[c(4, 3)]],
            fun = savi_fun)
#plot(savi, col = topo.colors(30), font.main = 2, main = "Soil Adjusted Vegetation Index (SAVI) for Landsat-9 OLI/TIRS C1 image \nLC09_L1TP_197055_20220111_20220122_02_T1: Yamassoukro, central C\^ote d'Ivoire", cex.main=1.0)
plot(savi, col = brewer.pal(11, "Accent"), font.main = 2, main = "Soil Adjusted Vegetation Index (SAVI) for Landsat-9 OLI/TIRS C1 image \nLC09_L1TP_197055_20220111_20220122_02_T1: Yamassoukro, central C\^ote d'Ivoire", cex.main=1.0)
# Computing histogram for SAVI
hist(savi, font.main = 1, main = "SAVI values for Landsat-9 OLI/TIRS C1 image \n LC09_L1TP_197055_20220111_20220122_02_T1, Yamassoukro, central C\^ote d'Ivoire", xlab = "SAVI", ylab= "Frequency",
    col = "darkgoldenrod1", xlim = c(-0.4, 0.4),  breaks = 30, xaxt = "n")
axis(side=1, at = seq(-0.4, 0.4, 0.1), labels = seq(-0.4, 0.4, 0.1))
#
# 3. Enhanced Vegetation Index 2
# EVI2 = 2.4 * (Band 5 - Band 4) / (Band 5 + Band 4 + 1).
evi2_fun = function(nir, red){
  2.4 * ((nir - red) / (nir + red + 1))
}
evi2 = lapp(landsat[[c(4, 3)]],
            fun = evi2_fun)
plot(evi2, col = brewer.pal(9, "Set1"), font.main = 2, main = "Enhanced Vegetation Index 2 (EVI2) for Landsat-9 OLI/TIRS C1 image \nLC09_L1TP_197055_20220111_20220122_02_T1: Yamassoukro, central C\^ote d'Ivoire", cex.main=0.9)
# Computing histogram for EVI2
hist(evi2, font.main = 1, main = "EVI2 values for Landsat-9 OLI/TIRS C1 image \n LC09_L1TP_197055_20220111_20220122_02_T1, Yamassoukro, central C\^ote d'Ivoire", xlab = "SAVI", ylab= "Frequency",
    col = "darkorchid1", xlim = c(-0.4, 0.4),  breaks = 30, xaxt = "n")
axis(side=1, at = seq(-0.4, 0.4, 0.1), labels = seq(-0.4, 0.4, 0.1))
#
# 4. Atmospherically Resistant Vegetation Index 2 (ARVI2)
# ARVI2 = -0.18 + 1.17 ((nir - red) / (nir + red))
arvi_fun = function(nir, red){
  ((nir - red) / (nir + red)) * 1.17 - 0.18
}
arvi = lapp(landsat[[c(4, 3)]],
            fun = arvi_fun)
plot(arvi, col = brewer.pal(12, "Paired"), font.main = 1, main = "Atmospherically Resistant Vegetation Index 2 (ARVI2) for Landsat-9 OLI/TIRS C1 \nimage LC09_L1TP_197055_20220111_20220122_02_T1: Yamassoukro, C\^ote d'Ivoire", cex.main=0.9)
# Computing histogram for ARVI2
hist(arvi, font.main = 1, main = "ARVI2 values for Landsat-9 OLI/TIRS C1 image \n LC09_L1TP_197055_20220111_20220122_02_T1, Yamassoukro, C\^ote d'Ivoire", xlab = "ARVI2", ylab= "Frequency",
    col = "cyan1", xlim = c(-0.4, 0.1),  breaks = 30, xaxt = "n")
axis(side=1, at = seq(-0.4, 0.1, 0.1), labels = seq(-0.4, 0.1, 0.1))
\end{lstlisting}

\subsection{Terrain Analysis in Python}
The problem of terrain analysis in image processing can be referred as spatial data modelling with classification of cells according to values of elevations and obtaining relevant derivatives and variables for visualization of slope, aspect and hillshade using parameters of sun and azimuth. The 3D visualization of the terrain surface is based on modelling azimuth and altitude of the artificial light imitating the sun illumination.

The terrain analysis has been performed using integration of Python libraries EarthPy \cite{wasser2019EarthPyAP}, Matplotlib \cite{Hunter2007MatplotlibA2}, rasterio \cite{gillies_2019}, and auxiliary packages NumPy \cite{harris2020array} and Plotly \cite{plotly}.

\begin{lstlisting}[language=python, caption=Python code (1) used for plotting DEM in Figure 2, label={lst:08}]
python3.10
# Set home directory
import os
os.chdir('/Users/polinalemenkova/Documents/Python/Image_Processing')
os. getcwd()
# Set data source
dtm = "n07_w006_3arc_v2.tif"
# Load libraries
from glob import glob
import earthpy as et
import earthpy.spatial as es
import earthpy.plot as ep
import rasterio as rio
import matplotlib.pyplot as plt
import numpy as np
from matplotlib.colors import ListedColormap
import plotly.graph_objects as go
np.seterr(divide='ignore', invalid='ignore')
# Open DEM by Rasterio which reads GeoTIFF format and stores gridded raster datasets of digital terrain models (DTMs) and satellite images
with rio.open(dtm) as src:
    elevation = src.read(1)
# Plot the data
fig, ax = plt.subplots(figsize=(10, 10))
ep.plot_bands(
    elevation,
    ax=ax,
    cmap="terrain",
    title="Ditigal Terrain Model (DTM) of Kossou Lake, C\^ote d'Ivoire. \n Data source: SRTM3 (3 arc-second resolution, 90 m), NASA/NGA. \nEntity ID: SRTM3N07W006V2. Aquisition date: 2000-02-11. cpt: 'terrain'",
    figsize=(10, 6),
)
plt.show()
# Save figure with 300 dpi resolution
plt.savefig('Figure_2_DEM_Kossou_1609.jpg', dpi=300)
\end{lstlisting}

\begin{lstlisting}[language=python, caption=Python code (2) used for plotting hillshade in Figure 3, label={lst:09}]
python3.10
# Set home directory
import os
os.chdir('/Users/polinalemenkova/Documents/Python/Image_Processing')
os. getcwd()
# Set data source
dtm = "n07_w006_3arc_v2.tif"
# Load libraries
from glob import glob
import earthpy as et
import earthpy.spatial as es
import earthpy.plot as ep
import rasterio as rio
import matplotlib.pyplot as plt
import numpy as np
from matplotlib.colors import ListedColormap
import plotly.graph_objects as go
np.seterr(divide='ignore', invalid='ignore')
# Generate hillshade from DTM
hillshade = es.hillshade(elevation)
# Plot the data
fig, ax = plt.subplots(figsize=(10, 10))
ep.plot_bands(
    hillshade,
    ax=ax,
    cbar=True,
    cmap="cividis",
    title="Hillshade from DTM of Kossou Lake, C\^ote d'Ivoire. \n Data source: SRTM3 (3 arc-second resolution, 90 m), NASA/NGA. \nEntity ID: SRTM3N07W006V2. Aquisition date: 2000-02-11. cmap=cividis",
    figsize=(10, 10),
)
plt.show()
plt.savefig('Figure_3_Hillshade_Kossou_1609.png', dpi=300)
\end{lstlisting}

\begin{lstlisting}[language=python, caption=Python code (3) and (4) used for plotting hillshade with changed azimuth and sun angle values for Figures 4 and 5, label={lst:10}]
python3.10
# Set home directory
import os
os.chdir('/Users/polinalemenkova/Documents/Python/Image_Processing')
os. getcwd()
# Set data source
dtm = "n07_w006_3arc_v2.tif"
# Load libraries
from glob import glob
import earthpy as et
import earthpy.spatial as es
import earthpy.plot as ep
import rasterio as rio
import matplotlib.pyplot as plt
import numpy as np
from matplotlib.colors import ListedColormap
import plotly.graph_objects as go
np.seterr(divide='ignore', invalid='ignore')
# Open DEM by Rasterio which reads GeoTIFF format and stores gridded raster datasets of digital terrain models (DTMs) and satellite images
with rio.open(dtm) as src:
    elevation = src.read(1)
# Changing the azimuth of the sun of the hillshade layer to 210 degrees
hillshade_azimuth_230 = es.hillshade(elevation, azimuth=230)
# Plotting the hillshade layer with the modified azimuth
fig, ax = plt.subplots(figsize=(10, 10))
ep.plot_bands(
    hillshade_azimuth_230,
    ax=ax,
    cbar=True,
    cmap="plasma",
    title="Hillshade from DTM of Kossou Lake, C\^ote d'Ivoire. Azimuth of sun is adjusted to 230 \N{DEGREE SIGN}",
    figsize=(10, 10),
)
plt.show()
plt.savefig('Figure_4_Azimuth_hillshade_Kossou.jpg', dpi=300)
# Changing the angle altitude of sun
hillshade_angle_10 = es.hillshade(elevation, altitude=10)
# Plotting the hillshade layer with the modified angle altitude
fig, ax = plt.subplots(figsize=(10, 10))
ep.plot_bands(
    hillshade_angle_10,
    ax=ax,
    cbar=True,
    cmap="magma",
    title="Hillshade from DTM of Kossou Lake, C\^ote d'Ivoire. Angle altitude is defined as 10 \N{DEGREE SIGN}",
    figsize=(10, 10),
)
plt.show()
plt.savefig('Figure_5_Angle_hillshade_Kossou.jpg', dpi=300)
\end{lstlisting}

%%%%%%%%%%%%%%%%%%%%%%%%%%%%%%%%%%%%%%%%%%
\newpage
\section{Results}
%--------
\subsection{Color composites}
In Figure \ref{fig06}c) Bands 3-4-5 were taken in inverse way where near-infrared band (B5) was representing red, red (B4) as green, and green (B3) as blue.

\begin{figure}[H]
\makebox[0.90\linewidth][r]{%
	\begin{subfigure}[b]{.6\textwidth}
		\centering
			\includegraphics[width=0.95\textwidth]{Fig_06a}
		\caption{Bands 2-3-4, stretch improvement: 'linear'}
	\end{subfigure}%
	\begin{subfigure}[b]{.6\textwidth}
		\centering
		\includegraphics[width=0.95\textwidth]{Fig_06b}
		\caption{Bands 2-3-4, stretch improvement: 'histogram'.}
	\end{subfigure}%
}\\
\vfill \vspace{1mm}
\makebox[0.90\linewidth][r]{%
	\begin{subfigure}[b]{.6\textwidth}
		\centering
			\includegraphics[width=0.95\textwidth]{Fig_06c}
		\caption{Bands 3-4-5.}
	\end{subfigure}%
	\begin{subfigure}[b]{.6\textwidth}
		\centering
		\includegraphics[width=0.95\textwidth]{Fig_06d}
		\caption{Bands 7-6-5}
	\end{subfigure}%
}
\caption{Natural color composites: (a) and (b). False color composites: (c) and (d). Landsat 9 OLI image of Kossou Lake region and Yamoussoukro, Côte d'Ivoire. Mapping: RStudio. Source: authors.}\label{fig06}
\end{figure}

\begin{figure}[H]
\makebox[0.90\linewidth][c]{%
	\begin{subfigure}[b]{.6\textwidth}
		\centering
			\includegraphics[width=0.95\textwidth]{Fig_07a}
		\caption{Bands 5-7-3.}
	\end{subfigure}%
	\begin{subfigure}[b]{.6\textwidth}
		\centering
			\includegraphics[width=0.95\textwidth]{Fig_07b}
		\caption{Bands 7-4-2}
	\end{subfigure}%
}\\
 \vfill \vspace{1mm}
\makebox[0.90\linewidth][c]{%
	\begin{subfigure}[b]{.6\textwidth}
		\centering
			\includegraphics[width=0.95\textwidth]{Fig_07c}
		\caption{Bands 5-4-3 (Color Infrared).}
	\end{subfigure}%
	\begin{subfigure}[b]{.6\textwidth}
		\centering
			\includegraphics[width=0.95\textwidth]{Fig_07d}
		\caption{Bands 6-5-4}
	\end{subfigure}%
}
\caption{False color composite visualization of Landsat 8 OLI image of Kossou Lake surroundings, Côte d'Ivoire. Mapping: RStudio. Source: authors.}\label{fig07}
\end{figure}

\subsection{Computing Vegetation Indices}
\subsubsection{Normalized Difference Vegetation Index (NDVI)}
\begin{figure}[H]
\makebox[0.97\linewidth][r]{%
	\begin{subfigure}[b]{.68\textwidth}
		\centering
			\includegraphics[width=0.95\textwidth]{Fig_08a.jpg}
		\caption{NDVI ratio based vegetation index}
	\end{subfigure}%
	\begin{subfigure}[b]{.61\textwidth}
		\centering
			\includegraphics[width=0.95\textwidth]{Fig_08b.jpg}
		\caption{Data distribution within the computed NDVI}
	\end{subfigure}%
}
       \caption{Normalized Difference Vegetation Index (NDVI) computed for multispectral image \emph{LC09\_L1TP\_197055\_20220111\_20220122\_02\_T1\_B}.}
        \label{fig08}
\end{figure}

\subsubsection{Soil Adjusted Vegetation Index (SAVI)}
\begin{figure}[H]
\makebox[0.97\linewidth][r]{%
	\begin{subfigure}[b]{.68\textwidth}
		\centering
			\includegraphics[width=0.95\textwidth]{Fig_09a.jpg}
		\caption{SAVI ratio based vegetation index}
	\end{subfigure}%
	\begin{subfigure}[b]{.61\textwidth}
		\centering
			\includegraphics[width=0.95\textwidth]{Fig_09b.jpg}
		\caption{Data distribution within the computed SAVI}
	\end{subfigure}%
}
       \caption{Soil Adjusted Vegetation Index (SAVI) computed for multispectral image \emph{LC09\_L1TP\_197055\_20220111\_20220122\_02\_T1\_B}.}
        \label{fig09}
\end{figure}

\subsubsection{Enhanced Vegetation Index 2 (EVI2)}
\begin{figure}[H]
\makebox[0.97\linewidth][r]{%
	\begin{subfigure}[b]{.68\textwidth}
		\centering
			\includegraphics[width=0.95\textwidth]{Fig_10a.jpg}
		\caption{EVI2 ratio based vegetation index}
	\end{subfigure}%
	\begin{subfigure}[b]{.61\textwidth}
		\centering
			\includegraphics[width=0.95\textwidth]{Fig_10b.jpg}
		\caption{Data distribution within the computed EVI2}
	\end{subfigure}%
}
       \caption{Enhanced Vegetation Index 2 (EVI2) computed for multispectral image \emph{LC09\_L1TP\_197055\_20220111\_20220122\_02\_T1\_B}.}
        \label{fig10}
\end{figure}

\subsubsection{Atmospherically Resistant Vegetation Index 2 (ARVI2)}
\begin{figure}[H]
\makebox[0.97\linewidth][r]{%
	\begin{subfigure}[b]{.68\textwidth}
		\centering
			\includegraphics[width=0.95\textwidth]{Fig_11a.png}
		\caption{ARVI2 ratio based vegetation index}
	\end{subfigure}%
	\begin{subfigure}[b]{.61\textwidth}
		\centering
			\includegraphics[width=0.95\textwidth]{Fig_11b.jpg}
		\caption{Data distribution within the computed EVI2}
	\end{subfigure}%
}
       \caption{Atmospherically Resistant Vegetation Index 2 (ARVI2) computed for multispectral image \emph{LC09\_L1TP\_197055\_20220111\_20220122\_02\_T1\_B}.}
        \label{fig11}
\end{figure}

\newpage
\subsection{Terrain Analysis}
\begin{figure}[H]
     \centering
     \begin{subfigure}[b]{0.45\textwidth}
         \centering
         \includegraphics[width=8.6cm]{Fig_12a}
         \caption{}
         \label{fig:00a}
     \end{subfigure} \hfill \hspace{2mm}
     \begin{subfigure}[b]{0.45\textwidth}
         \centering
         \includegraphics[width=7cm]{Fig_12b}
         \caption{}
         \label{fig:00a}
     \end{subfigure}
        \caption{\textbf{(a)}: Landsat OLI/TIRS image with enlarged fragment of Kossou Lake: grayscale visualization of SWIR Band 6. Mapping: R. \textbf{(b)}: Enlarged fragment of the cropped SRTM3 Void Filled DEM overlayed on top of the hillshade. cmaps: 'turbo' and 'gist\_yarg'. Mapping: Python.}
        \label{fig12}
\end{figure}

\begin{figure}[H]
\makebox[0.90\linewidth][c]{%
	\begin{subfigure}[b]{.6\textwidth}
		\centering
			\includegraphics[width=0.95\textwidth]{Fig_13a}
		\caption{DTM of Kossou Lake, Côte d'Ivoire.}
	\end{subfigure}%
	\begin{subfigure}[b]{.6\textwidth}
		\centering
			\includegraphics[width=0.95\textwidth]{Fig_13b}
		\caption{Hillshade from DTM of Kossou Lake, Côte d'Ivoire.}
	\end{subfigure}%
}\\
 \vfill \vspace{2mm}
\makebox[0.90\linewidth][c]{%
	\begin{subfigure}[b]{.6\textwidth}
		\centering
			\includegraphics[width=0.95\textwidth]{Fig_13c}
		\caption{Azimuth of sun: 230\textdegree , hillshade from DTM.}
	\end{subfigure}%
	\begin{subfigure}[b]{.6\textwidth}
		\centering
			\includegraphics[width=0.95\textwidth]{Fig_13d}
		\caption{Sun angle altitude: 10\textdegree , hillshade from DTM.}
	\end{subfigure}%
}
\caption{Terrain modelling for Kossou Lake area, Côte d'Ivoire. Mapping: Python. Source: authors.} \label{fig13}
\end{figure}

%%%%%%%%%%%%%%%%%%%%%%%%%%%%%%%%%%%%%%%%%%
\newpage
\section{Discussion}


%%%%%%%%%%%%%%%%%%%%%%%%%%%%%%%%%%%%%%%%%%
\section{Conclusions}
Integrated and wider use of regional environmental resources present the potential for sustainable development of Côte-d'Ivoire. This requires regular monitoring of the environmental conditions of the country using satellite images. In the data science and analysis, processing satellite images using advanced programming methods is a challenging problem. In this paper, we presented the solution of this problem by the straightforward yet effective approach based on combining libraries of Python and R languages for processing of raster data in a scripting framework. 

First, the terrain surface materials were differentiated by the RGB triplet colours by plotting the colour composite of the Landsat 98 OLI/TIRS image using R library 'raster'. The colours correspond to the composites of the RGB signals scattered from the natural and built environment which resulted in various DN values of pixel cells in the raster matrix. The tested combinations of the three channels as sensor bands enabled to display the image and visualise the appearance of various objects on the terrain using distinct colours using R library 'terra'. We compared the correlation between Landsat bands for selected Landsat-8 OLI/TIRS C1 bands in multispectral image for bands 1 and 2 (Blue and Green) and 4 and 5 (red and NOR), correspondingly.

We then computed the three selected vegetation indices - NDVI, SAVI and EVI2 using the principles of map algebra based on DNs of reflectance value of pixels in selected spectral bands -- NIR and red. As demonstrated, the major effect of spectral reflectance from healthy green vegetation in visible in the NIR region of the spectrum of the image with the strong NIR response from it raising sharply for  healthy vegetation in excellent conditions. Conversely, areas with poor or damaged vegetation are represented by low NIR response which decreases in the maps of the vegetation indices. 

Finally, the terrain modelling was performed using the SRTM DEM with high resolution using the algorithms of Python library EarthPy by combining the topographic 3D data with spatial information for geomorphic visualization of slope, aspect and elevation for the relief analysis. We presented the methods as a series of scripts with detailed explanation of the functionality of codes and effects on parameters of the image.

The main advantage of our study is that deriving the information from the remotely sensed data using programming languages presents the efficient and straightforward approach designed to handle complex scenes of the remotely sensed data. Since the environmental dynamics of the landscapes of the Earth usually involves long-term and short-term variations, this includes processing large sequences of data as batch series where conventional GIS can not process successfully and which requires automatic image processing. A possible way to handle this issue includes the advanced methods of high-level programming languages, such as Python and R, which can process large Earth datasets using libraries for geospatial data processing, as demonstrated in this paper.

Using multiple data processed by programming algorithms, it is possible to depict the nonlinear environmental changes and highlight complex relationships among the processes and objects in the landscapes of the Earth. These include a wide variety of phenomenas and processes that can be depicted on the maps using interpretation of the satellite images, e.g., climate change, urban sprawl, vegetation degradation, deforestation, geomorphic fluctuations in the relief that might be caused both by the geologic effects and anthropogenic activities.

%%%%%%%%%%%%%%%%%%%%%%%%%%%%%%%%%%%%%%%%%%

\funding{This research received no external funding.}

\institutionalreview{Not applicable}

\informedconsent{Not applicable}

\dataavailability{Not applicable} 

\acknowledgments{We thank the two anonymous reviewers for comments and reading of this manuscript}

\conflictsofinterest{The authors declare no conflict of interest} 

\abbreviations{Abbreviations}{
The following abbreviations are used in this manuscript:\\

\noindent 
\begin{tabular}{@{}ll}
CPT & Colour Palette Table \\
DCW & Digital Chart of the World \\
DN & Digital Number \\
DEM & Digital Elevation Model \\
DTM & Digital Terrain Model \\
EVI & Enhanced Vegetation Index \\
GEBCO & General Bathymetric Chart of the Oceans \\
GDAL & Geospatial Data Abstraction Library \\
GIS & Geographic Information System \\
GloVis & Global Visualization Viewer \\
GMT & Generic Mapping Tools \\
GPS & Global Positioning System \\
Landsat 8-9 OLI/TIRS & Landsat 8-9 Operational Land Imager and Thermal Infrared \\
LiDAR & Light Detection and Ranging \\ 
NASA & National Aeronautics and Space Administration \\
NDVI & Normalized Difference Vegetation Index \\
NGA & National Geospatial-Intelligence Agency \\
NIR & Near Infrared \\
SAVI & Soil-Adjusted Vegetation Index \\
SRTM & Shuttle Radar Topography Mission \\
SWIR & Short Wave InfraRed \\
USGS & United States Geological Survey \\
\end{tabular}
}

\reftitle{References}

% Please provide either the correct journal abbreviation (e.g. according to the “List of Title Word Abbreviations” http://www.issn.org/services/online-services/access-to-the-ltwa/) or the full name of the journal.
% Citations and References in Supplementary files are permitted provided that they also appear in the reference list here. 

%=====================================
% References, variant A: external bibliography
%=====================================
\bibliography{IvoryCoast.bib}
%\nocite{*}

%%%%%%%%%%%%%%%%%%%%%%%%%%%%%%%%%%%%%%%%%%
%\end{adjustwidth}
\end{document}

